\documentclass[../DoAn.tex]{subfiles}
\begin{document}

\section{Đặt vấn đề}
\label{section:1.1}

Trong bối cảnh tiến trình chuyển đổi số đang diễn ra mạnh mẽ và sâu rộng tại Việt Nam, nhu cầu tự động hóa các quy trình nghiệp vụ hành chính cơ bản trong các doanh nghiệp, cơ quan và tổ chức ngày càng trở nên cấp thiết. Một trong những nghiệp vụ hành chính cốt lõi nhưng đang tiêu tốn nhiều chi phí vận hành cũng như nguồn nhân lực quản lý trực tiếp nhất hiện nay chính là công tác quản lý chấm công, điểm danh và kiểm soát dòng người ra vào. 

Khảo sát thực tế tại các cơ quan, tổ chức cho thấy, bài toán kiểm soát an ninh ra vào hiện nay đang phải đối mặt với ba thách thức lớn chưa có lời giải triệt để. Thứ nhất là nguy cơ gian lận danh tính khi các phương thức truyền thống như ký tên trên giấy, nhập mã định danh hoặc thậm chí là các hệ thống thẻ từ nội bộ thông thường rất dễ bị chuyển nhượng, dẫn đến tình trạng cho mượn thẻ để chấm công hộ, làm sai lệch dữ liệu nhân sự. Thứ hai là gánh nặng logistics và chi phí quản lý khi tổ chức phải tự bỏ ngân sách để mua sắm, in ấn, phân phát, thu hồi và cấp lại các loại thẻ định danh riêng của đơn vị mỗi khi nhân sự làm mất hoặc hỏng hóc, tạo ra một quy trình hành chính cồng kềnh và tiêu tốn chi phí định kỳ. Thứ ba là sự thiếu hụt năng lực giám sát tập trung và phản hồi theo thời gian thực, khiến người quản lý hoàn toàn bị động, không thể bao quát trạng thái an ninh của toàn bộ các điểm kiểm soát trong hạ tầng đơn vị.

Nếu bài toán tối ưu hóa quy trình kiểm soát ra vào và tự động hóa điểm danh dựa trên các phương thức định danh chính thống sẵn có được giải quyết một cách triệt để, nó sẽ mang lại lợi ích to lớn và đa chiều cho toàn bộ hệ thống hành chính. Đối với tổ chức và doanh nghiệp, giải pháp này giúp cắt giảm tối đa chi phí cấp phát thẻ nội bộ, loại bỏ hoàn toàn các sai số và gian lận nhờ cơ chế đối soát đa yếu tố, từ đó tối ưu hóa hiệu suất lao động. Đối với nhân sự, hệ thống sẽ tận dụng ngay giấy tờ định danh chính thức do Nhà nước cấp để thực hiện quyền truy cập mà không cần phải mang theo nhiều loại thẻ hành chính trung gian khác nhau. Không dừng lại ở phạm vi văn phòng công sở, một mô hình quản lý ra vào an toàn, bảo mật cao và minh bạch như vậy còn sở hữu tiềm năng ứng dụng rộng mở sang rất nhiều lĩnh vực khác của đời sống xã hội bao gồm quản lý cư dân tại các khu đô thị thông minh hoặc quản lý học viên tại các cơ sở giáo dục quy mô lớn.

Xuất phát từ những đòi hỏi thực tiễn mang tính cấp thiết và có quy mô ứng dụng rộng lớn đó, bài toán nghiên cứu thiết kế một hạ tầng kiểm soát ra vào thông minh, an toàn, sử dụng giải pháp xác thực đa yếu tố và quản lý tập trung chính là hướng đi trọng tâm được lựa chọn để thực hiện trong đồ án tốt nghiệp này.

\section{Mục tiêu và phạm vi đề tài}
\label{section:1.2}

Các sản phẩm kiểm soát ra vào và điểm danh tự động hiện nay trên thị trường tương đối đa dạng nhưng thường vận hành độc lập theo hai hướng tiếp cận riêng biệt bao gồm nhóm giải pháp dựa trên các thiết bị quét sinh trắc học cố định như máy quét vân tay hoặc camera nhận diện khuôn mặt chuyên dụng, và nhóm giải pháp dựa trên các thiết bị đọc thẻ vật lý cục bộ như thẻ từ RFID hoặc thẻ thông minh nội bộ của tổ chức. Thông qua quá trình khảo sát và đánh giá tổng quan, các hệ thống hiện tại vẫn tồn tại ba hạn chế lớn. Thứ nhất, tính bảo mật chưa cao khi các hệ thống thẻ từ thông thường không thể xác thực được người đang cầm thẻ có phải là chính chủ hay không, trong khi các hệ thống nhận diện khuôn mặt thương mại phân khúc phổ thông lại dễ bị đánh lừa bằng các kỹ thuật giả mạo đơn giản như sử dụng ảnh in hoặc video phát lại. Thứ hai, gánh nặng đăng ký và lưu trữ sinh trắc học khi tổ chức bắt buộc phải tiến hành thu thập, đăng ký ảnh mẫu của từng nhân sự một cách thủ công, đồng thời phải lưu trữ toàn bộ kho ảnh gốc tập trung trên máy chủ, gây ra các rủi ro lớn về an toàn thông tin và vi phạm quyền riêng tư dữ liệu cá nhân. Thứ ba là sự thiếu đồng bộ và chi phí triển khai lớn do các sản phẩm kiểm soát an ninh chuyên dụng thường sử dụng các giao thức đóng, đòi hỏi đầu tư hạ tầng phần cứng đồng bộ đắt đỏ và phức tạp trong việc tích hợp với các phần mềm quản lý hành chính sẵn có của đơn vị.

Nhằm khắc phục triệt để các hạn chế nêu trên, đề tài hướng tới giải quyết bài toán kiểm soát an ninh thông qua việc tận dụng cơ sở dữ liệu quốc gia sẵn có trên thẻ Căn cước công dân và thiết lập cơ chế xác thực kép. Trên cơ sở đó, mục tiêu tổng quát của đề tài là thiết kế, cài đặt và thử nghiệm một hệ thống điểm danh và kiểm soát ra vào hoàn chỉnh, tích hợp nhận diện khuôn mặt và bóc tách dữ liệu gốc mã hóa từ chip bảo mật NFC của thẻ Căn cước công dân, hướng đến khả năng triển khai thực tế tối ưu trong môi trường doanh nghiệp hoặc các cơ sở giáo dục.

Để hiện thực hóa mục tiêu tổng quát, các mục tiêu cụ thể của đề tài bao gồm năm nhiệm vụ trọng tâm được thực hiện tuần tự : (i) xây dựng ứng dụng trạm đầu cuối di động trên nền tảng Flutter có khả năng bắt tay bảo mật giao tiếp trường gần NFC để bóc tách thông tin chân dung từ chip Căn cước công dân kết hợp chụp ảnh trực tiếp tại chỗ; (ii) triển khai dịch vụ trí tuệ nhân tạo trên máy chủ chịu trách nhiệm trích xuất vector đặc trưng khuôn mặt và đối sánh không gian với độ chính xác cao; (iii) xây dựng hệ thống điều phối Backend trung tâm dựa trên kiến trúc REST API quản lý toàn bộ luồng dữ liệu, xử lý logic xác thực hai lớp và kiểm tra ma trận phân quyền; (iv) phát triển firmware điều khiển thiết bị IoT ngoại vi cho vi điều khiển ESP32 nhận lệnh qua giao thức mạng để đóng ngắt rơ-le vật lý điều khiển khóa cửa ; và (v) cung cấp giao diện bảng quản trị trực tuyến vận hành trên cổng mạng 3002 hiển thị dữ liệu dòng người ra vào theo thời gian thực, tự động phân tách và làm nổi bật các trạng thái cảnh báo chuỗi điểm danh phục vụ hậu kiểm, kết xuất báo cáo dữ liệu định kỳ.

Tính đột phá mang lại giá trị khoa học và thực tiễn của đề tài đồ án này nằm ở giải pháp kiến trúc xác thực hai lớp không cần đăng ký trước sinh trắc học. Hệ thống không yêu cầu tổ chức phải thu thập hay lưu trữ kho ảnh gốc của nhân viên lên máy chủ từ trước để làm ảnh đối chiếu. Thay vào đó, mỗi khi nhân sự thực hiện điểm danh, hệ thống sử dụng chính tệp tin ảnh chân dung gốc chất lượng cao được bảo chứng bởi chữ ký số của cơ quan có thẩm quyền bên trong chip thẻ Căn cước công dân làm dữ liệu tham chiếu tức thời tại chỗ. Giải pháp này vừa nâng mức độ an toàn an ninh lên tối đa, vừa bảo vệ tuyệt đối quyền riêng tư dữ liệu sinh trắc học của người dùng, đồng thời tinh giản hoàn toàn quy trình vận hành hành chính của tổ chức.

\section{Định hướng giải pháp}
\label{section:1.3}

Để giải quyết triệt để các nhiệm vụ khoa học và kỹ thuật đã được xác định, đồ án đề xuất định hướng giải pháp tổng thể dựa trên sự phân tách kiến trúc và tích hợp đa nền tảng, được định hình xây dựng dựa trên sự liên kết của bốn nhóm giải pháp công nghệ nền tảng. 

Thứ nhất là công nghệ nhận diện khuôn mặt dựa trên học sâu ứng dụng thuật toán mạng nơ-ron tích chập sâu ArcFace~\cite{deng2019arcface} thông qua thư viện InsightFace~\cite{insightface_github} kết hợp cùng mô hình FaceNet~\cite{facenet2015} làm phương án dự phòng, giúp bóc tách hình học khuôn mặt thành các vector đặc trưng trong không gian embedding và đối sánh qua khoảng cách Cosine chính xác. Thứ hai là công nghệ bóc tách dữ liệu chip bảo mật NFC theo giao thức bắt tay an toàn BAC/PACE tuân thủ tiêu chuẩn quốc tế ICAO Doc 9303~\cite{icao9303}, thực thi thông qua thư viện hỗ trợ trên môi trường Flutter nhằm bóc tách hợp pháp phân vùng dữ liệu chứa ảnh chân dung gốc bên trong chip thẻ. Thứ ba là kiến trúc hệ thống phân tán hướng dịch vụ phân rã thành các module độc lập bao gồm lõi xử lý nghiệp vụ Backend, dịch vụ tính toán trí tuệ nhân tạo chuyên biệt và phân hệ hiển thị người dùng nhằm tối ưu hóa tải tài nguyên phần cứng. Thứ tư là giao thức truyền thông mạng thời gian thực tích hợp song song giao thức nhắn tin gọn nhẹ MQTT~\cite{mqtt5} cho kênh điều khiển thiết bị phần cứng ngoại vi và giao thức STOMP qua WebSocket~\cite{rfc6455} cho kênh đồng bộ dữ liệu giao diện web hiển thị nhật ký tức thời.

Dựa trên các định hướng công nghệ nêu trên, đồ án xây dựng một giải pháp phần mềm kiểm soát ra vào đồng bộ kết hợp thiết bị di động thông minh, máy chủ và mạch chấp hành IoT. Khi nhân sự thực hiện quét thẻ, ứng dụng trạm đầu cuối Flutter thực hiện giao tiếp đọc tệp tin dữ liệu từ chip Căn cước công dân qua kết nối NFC để làm ảnh mẫu tham chiếu gốc, đồng thời kích hoạt camera chụp ảnh thực tế để thực hiện tiến trình đối soát lớp thứ nhất ngay tại chỗ. Gói dữ liệu sau đó được đẩy về máy chủ Backend để kích hoạt lớp xác thực thứ hai chuyên sâu tại trung tâm thông qua việc gọi dịch vụ trí tuệ nhân tạo đối sánh chéo vector đặc trưng khuôn mặt và kiểm tra ma trận phân quyền truy cập. Khi các lớp xác thực danh tính sinh trắc học cho kết quả hợp lệ, máy chủ Backend sẽ thực hiện đối soát chuỗi trạng thái di chuyển thời gian thực. Trong trường hợp nhân sự quét sai chuỗi logic tuần tự (thiếu lượt quẹt VÀO hoặc thiếu lượt quẹt RA trong cùng một chu kỳ làm việc) nhưng đảm bảo thông tin chính chủ (Face Match đạt tiêu chuẩn an ninh), hệ thống tự động kích hoạt kịch bản an ninh mềm (Soft Mode). Lúc này, bộ điều phối trung tâm sẽ đồng thời thực hiện các tác vụ thời gian thực bao gồm phát lệnh mở khóa thông qua thông điệp MQTT xuống bo mạch vi điều khiển ESP32 để giải phóng chốt cửa vật lý, truyền thông điệp điều khiển hiển thị banner nhắc nhở trực quan trực tiếp trên màn hình trạm, và kích hoạt kênh truyền WebSocket STOMP để tự động đẩy dòng log mang nhãn cảnh báo chuyên biệt (như MISSING\_IN\_AUTO\_OPEN hoặc ANTI\_PASSBACK\_IN) hiển thị dạng Badge màu cam trực tuyến trên màn hình hiển thị Web Dashboard của người quản lý.

Đóng góp khoa học và thực tiễn lớn nhất của đồ án tốt nghiệp này là thiết kế và chứng minh thành công tính khả thi của một mô hình kiến trúc điểm danh hai lớp tự trị vận hành hoàn toàn dựa trên hạ tầng định danh quốc gia sẵn có. Kết quả đạt được của đề tài là một hệ thống phần mềm phân tán hoàn chỉnh, chạy ổn định, có khả năng ứng dụng trực tiếp vào môi trường doanh nghiệp quy mô thực tế, giúp loại bỏ công tác tổ chức thu thập dữ liệu mẫu sinh trắc học thủ công cho nhân sự mới, đồng thời loại bỏ hoàn toàn gánh nặng rủi ro pháp lý cũng như kỹ thuật trong việc xây dựng, bảo mật kho lưu trữ cơ sở dữ liệu ảnh khuôn mặt gốc tập trung trên hệ thống máy chủ của tổ chức.

\section{Bố cục đồ án}
\label{section:1.4}

Phần còn lại của báo cáo đồ án tốt nghiệp này được tổ chức thành năm chương tiếp theo với nội dung được phân bổ chi tiết.

Chương 2 trình bày kết quả khảo sát thực tế và phân tích yêu cầu hệ thống. Nội dung chương tập trung xác định chi tiết các tác nhân, yêu cầu chức năng và phi chức năng thông qua quá trình khảo sát trực tiếp tại doanh nghiệp, sau đó biểu diễn trực quan các yêu cầu dưới dạng biểu đồ Use Case tổng quan, các biểu đồ tuần tự mô tả luồng nghiệp vụ chính, biểu đồ lớp cấu trúc và biểu đồ triển khai vật lý hạ tầng, đồng thời thiết kế mô hình cơ sở dữ liệu quan hệ và đặc tả hệ thống giao diện lập trình API.

Chương 3 cung cấp nền tảng lý thuyết chuyên sâu và mô tả chi tiết các phân hệ công nghệ được sử dụng trong đồ án. Nội dung chương tập trung làm rõ nguyên lý học sâu nhận diện khuôn mặt theo hướng tiếp cận định lượng học máy với hàm mất mát ArcFace và kiến trúc FaceNet, tiêu chuẩn quốc tế ICAO Doc 9303 quy định cấu trúc phân vùng dữ liệu chip NFC trên thẻ Căn cước công dân, kết hợp các giao thức truyền thông đồng bộ và bất đồng bộ mạng bao gồm MQTT, REST API, kiến trúc mã hóa khóa bảo mật JWT kết hợp phân quyền RBAC, hệ thống framework Spring Boot, React và nền tảng di động Flutter.

Chương 4 trình bày chi tiết kết quả triển khai cài đặt hạ tầng công nghệ và hiện thực hóa mã nguồn hệ thống. Chương này tập trung mô tả quy trình xây dựng dịch vụ máy chủ trí tuệ nhân tạo Python Flask, cấu hình nghiệp vụ xác thực hai lớp không lưu trữ ảnh mẫu tại trung tâm trên nền tảng Java Spring Boot, lập trình giao diện Dashboard React cập nhật dữ liệu trực tuyến trên cổng 3002, và phát triển mã nguồn nhúng cho vi điều khiển ESP32 tiếp nhận lệnh mạng đóng ngắt rơ-le vật lý.

Chương 5 trình bày các kết quả đánh giá thực nghiệm, kịch bản kiểm thử hệ thống và các đóng góp kỹ thuật nổi bật của đề tài. Nội dung bao gồm quy trình kiểm thử đơn vị tích hợp các điểm cuối dịch vụ API, kiểm thử thực nghiệm đo lường sai số và độ chính xác của mô hình nhận diện khuôn mặt tại các ngưỡng khoảng cách hình học khác nhau, kiểm thử hiệu năng đầu cuối của toàn bộ luồng điểm danh vật lý bao gồm cả kịch bản kiểm thử tính năng Soft Mode hai chiều, kết hợp phân tích đánh giá an toàn bảo mật hệ thống.

Chương 6 tổng kết toàn bộ nội dung nghiên cứu và phát triển của đồ án tốt nghiệp. Chương này tiến hành đối chiếu cụ thể các kết quả thực tế đạt được so với mục tiêu đặt ra ban đầu, thẳng thắn chỉ ra những giới hạn kỹ thuật hiện tại của hệ thống, từ đó đề xuất các định hướng nghiên cứu và phát triển công nghệ khả thi trong tương lai nhằm tối ưu hóa hiệu năng và mở rộng quy mô ứng dụng thực tế.

\end{document}